\documentclass[12pt,a4paper]{article}
\usepackage[utf8]{inputenc}
\usepackage[T2A]{fontenc}
\usepackage[russian]{babel}
\usepackage{amsmath,amssymb}
\usepackage[left=25mm,right=25mm,top=20mm,bottom=20mm]{geometry}
\usepackage{enumitem}

\newcommand{\dd}{\mathrm{d}}
\newcommand{\pp}{\partial}
\newcommand{\ext}{\mathrm{ext}}
\newcommand{\eff}{\mathrm{eff}}
\newcommand{\st}{\text{ст}}
\newcommand{\sat}{\mathrm{sat}}

\title{Метод интегральных соотношений (моментов)\\
для задачи двухфазной фильтрации Black~Oil\\
с балансом энергии в радиальной геометрии}
\author{}
\date{}

\begin{document}
\maketitle
\tableofcontents
\newpage

%=============================================================================
\section{Введение и мотивация}\label{sec:intro}
%=============================================================================

Полная система уравнений двухфазной фильтрации Black~Oil с балансом
энергии выведена в~\cite{draft4}. Система состоит из трёх PDE
на $p(r,t)$, $T(r,t)$, $S_g(r,t)$ в области $r\in[r_w, R_\ext]$,
$t \geqslant 0$.

Стратегия приближённого решения: метод интегральных соотношений
(моментов) Баренблатта~\cite{barenblatt1954} применяется ко
\emph{всем трём уравнениям}. PDE сводятся к системе ODE для
параметров профилей $p(r,t)$, $S_g(r,t)$ и $T(r,t)$.
Температурное поле входит в уравнения давления и
насыщенности через PVT-свойства $R_s(p,T)$, $B_o(p,T)$,
$\mu_o(p,T)$ и т.\,д., а также через
$C_T\,\pp T/\pp t$ и $\beta_o^\sat\,\pp T/\pp t$,
поэтому <<последовательное>> решение
(сначала гидродинамика, потом температура) приводит
к неконтролируемым ошибкам.


%=============================================================================
\section{Исходные уравнения}\label{sec:equations}
%=============================================================================

Система состоит из трёх уравнений на три неизвестных
$p(r,t)$, $T(r,t)$, $S_g(r,t)$. Все PVT-свойства зависят
от $p$ и $T$: $R_s(p,T)$, $B_o(p,T)$, $B_g(p,T)$,
$\mu_\alpha(p,T)$ и т.\,д.

\subsection{Уравнение давления}

Из~\cite{draft4}, уравнение~(pressure):
\begin{equation}\label{eq:p_eq}
  \boxed{
    C_p\,\frac{\pp p}{\pp t}
    + C_T\,\frac{\pp T}{\pp t}
    + \mathcal{D} = 0,
  }
\end{equation}
где:

\paragraph{Подвижности фаз:}
\begin{equation}\label{eq:mobility}
  \lambda_\alpha = \frac{k\,k_{r\alpha}(S_g)}{\mu_\alpha(p,T)},
  \qquad \alpha \in \{o, g\}.
\end{equation}

\paragraph{Суммарная приведённая подвижность:}
\begin{equation}\label{eq:Lambda}
  \Lambda(p,T,S_g) = \frac{\lambda_o + \lambda_g}{B_g}
  = \frac{k}{B_g(p,T)}\!\left(
      \frac{k_{ro}(S_g)}{\mu_o(p,T)}
      + \frac{k_{rg}(S_g)}{\mu_g(p,T)}
    \right).
\end{equation}

\paragraph{Коэффициент ёмкости при $\pp p/\pp t$:}
\begin{equation}\label{eq:Cp}
  C_p = \phi\!\left[
    S_o\,(c_r + c_o^\sat)\!\left(\frac{1}{B_g} - \frac{R_s}{B_o}\right)
    + \xi_o\!\left(\frac{\pp R_s}{\pp p}
      + R_s\,(c_r + c_o^\sat)\right)
    + \frac{S_g}{B_g}\,(c_g + c_r)
  \right],
\end{equation}
где $\xi_o = S_o/B_o$, $S_o = 1 - S_g$, а <<полная>> сжимаемость
нефтяной фазы вдоль кривой насыщения:
\begin{equation}\label{eq:co_sat}
  c_o^\sat = c_o
    - \frac{1}{B_o}\frac{\pp B_o}{\pp R_s}\bigg|_{p,T}
      \frac{\pp R_s}{\pp p}\bigg|_T,
  \qquad
  c_o = -\frac{1}{B_o}\frac{\pp B_o}{\pp p}\bigg|_{T,\,R_s}.
\end{equation}

\paragraph{Коэффициент при $\pp T/\pp t$:}
\begin{equation}\label{eq:CT}
  C_T = \phi\!\left[
    -S_o\,\beta_o^\sat\!\left(\frac{1}{B_g} - \frac{R_s}{B_o}\right)
    + \xi_o\!\left(\frac{\pp R_s}{\pp T}
      - R_s\,\beta_o^\sat\right)
    - \frac{S_g}{B_g}\,\beta_g
  \right],
\end{equation}
где
\begin{equation}\label{eq:beta_o_sat}
  \beta_o^\sat = \beta_o
    + \frac{1}{B_o}\frac{\pp B_o}{\pp R_s}\bigg|_{p,T}
      \frac{\pp R_s}{\pp T}\bigg|_p,
  \qquad
  \beta_o = \frac{1}{B_o}\frac{\pp B_o}{\pp T}\bigg|_{p,\,R_s}.
\end{equation}

\paragraph{Пространственный оператор:}
\begin{equation}\label{eq:D_operator}
  \mathcal{D} = -\frac{1}{r}\,\frac{\pp}{\pp r}\!\left(
      r\,\Lambda\,\frac{\pp p}{\pp r}
    \right)
    - \frac{\lambda_o}{B_o}\!\left(
      \frac{\pp R_s}{\pp p}\,\frac{\pp p}{\pp r}
      + \frac{\pp R_s}{\pp T}\,\frac{\pp T}{\pp r}
    \right)\frac{\pp p}{\pp r}.
\end{equation}

\paragraph{Замечание.}
Знак при $\mathcal{D}$ в~\eqref{eq:p_eq}: $C_p\,\dot{p} + C_T\,\dot{T}
+ \mathcal{D} = 0$, т.\,е.\ $\mathcal{D}$ включает все
пространственные члены, перенесённые в левую часть.
В дивергентном слагаемом стоит минус, поскольку отток жидкости
($\mathcal{D} > 0$) снижает давление ($\dot{p} < 0$).


\subsection{Уравнение насыщенности}

Из~\cite{draft4}, уравнение~(Sg\_ode):
\begin{equation}\label{eq:Sg_eq}
  \boxed{
    \frac{\pp S_g}{\pp t}
    = S_o\,(c_r + c_o^\sat)\,\frac{\pp p}{\pp t}
      - S_o\,\beta_o^\sat\,\frac{\pp T}{\pp t}
      + \frac{B_o}{\phi}\,\frac{1}{r}\,\frac{\pp}{\pp r}\!
        \left(r\,\frac{u_o}{B_o}\right).
  }
\end{equation}
Подставляя $u_o = -\lambda_o\,\pp p/\pp r$:
\begin{equation}\label{eq:Sg_eq2}
  \frac{\pp S_g}{\pp t}
  = S_o\,(c_r + c_o^\sat)\,\frac{\pp p}{\pp t}
    - S_o\,\beta_o^\sat\,\frac{\pp T}{\pp t}
    - \frac{B_o}{\phi}\,\frac{1}{r}\,\frac{\pp}{\pp r}\!
      \left(\frac{r\,\lambda_o}{B_o}\,\frac{\pp p}{\pp r}\right).
\end{equation}

Член $-S_o\,\beta_o^\sat\,\pp T/\pp t$ описывает влияние
температуры на насыщенность: при охлаждении ($\pp T/\pp t < 0$)
растворимость газа растёт ($\pp R_s/\pp T < 0$ для типичных нефтей),
газ частично растворяется, и $S_g$ снижается. Знак зависит
от конкурирующих вкладов в $\beta_o^\sat$.


\subsection{Уравнение энергии}

Из~\cite{draft4}, уравнение~(energy\_final):
\begin{equation}\label{eq:energy}
  \boxed{
    (\rho c_p)_\eff\,\frac{\pp T}{\pp t}
    + (\rho c_p u)_\mathrm{conv}\,\frac{\pp T}{\pp r}
    = \frac{1}{r}\,\frac{\pp}{\pp r}\!\left(
        r\,\lambda_\eff\,\frac{\pp T}{\pp r}
      \right)
    + Q_\mathrm{JT} + Q_\mathrm{ad} + Q_\mathrm{deg},
  }
\end{equation}
где:

\paragraph{Эффективная теплоёмкость:}
\begin{equation}\label{eq:rhocp_eff}
  (\rho c_p)_\eff
  = (1-\phi)\,(\rho c_p)_s
    + \phi\,(S_o\,\rho_o\,c_{p,o} + S_g\,\rho_g\,c_{p,g}).
\end{equation}

\paragraph{Конвективный коэффициент:}
\begin{equation}\label{eq:rhocp_conv}
  (\rho c_p u)_\mathrm{conv}
  = \rho_o\,c_{p,o}\,u_o + \rho_g\,c_{p,g}\,u_g.
\end{equation}

\paragraph{Эффект Джоуля--Томсона:}
\begin{equation}\label{eq:QJT}
  Q_\mathrm{JT} = -\sum_\alpha u_\alpha\,(1 - \alpha_\alpha T)\,
    \frac{\pp p}{\pp r},
\end{equation}
где $\alpha_\alpha$ --- коэффициент теплового расширения
фазы~$\alpha$ при фиксированном составе.

\paragraph{Адиабатический эффект:}
\begin{equation}\label{eq:Qad}
  Q_\mathrm{ad} = \phi\!\left[
    T\sum_\alpha S_\alpha\,\alpha_\alpha
    + c_r\!\left(p + (\rho c_p)_s\,T\right)
  \right]\frac{\pp p}{\pp t}.
\end{equation}

\paragraph{Теплота дегазации:}
\begin{equation}\label{eq:Qdeg}
  Q_\mathrm{deg} = -L(p,T)\,\dot{m},
\end{equation}
где $L(p,T)$ --- удельная теплота дегазации, а массовый темп
выделения газа:
\begin{equation}\label{eq:mdot}
  \dot{m} = -\rho_g^\st\!\left[
    \phi\,\xi_o\!\left(
      \frac{\pp R_s}{\pp p}\,\frac{\pp p}{\pp t}
      + \frac{\pp R_s}{\pp T}\,\frac{\pp T}{\pp t}
    \right)
    + \frac{u_o}{B_o}\!\left(
      \frac{\pp R_s}{\pp p}\,\frac{\pp p}{\pp r}
      + \frac{\pp R_s}{\pp T}\,\frac{\pp T}{\pp r}
    \right)
  \right].
\end{equation}

\paragraph{Замечание.}
$Q_\mathrm{deg}$ содержит $\pp T/\pp t$ и $\pp T/\pp r$
через~\eqref{eq:mdot}. При подстановке в~\eqref{eq:energy}
они перегруппировываются с коэффициентами при $\pp T/\pp t$
и $\pp T/\pp r$:
\begin{equation}\label{eq:rhocp_eff_mod}
  (\rho c_p)_\eff^* = (\rho c_p)_\eff
    - L\,\rho_g^\st\,\phi\,\xi_o\,\frac{\pp R_s}{\pp T},
\end{equation}
\begin{equation}\label{eq:rhocp_conv_mod}
  (\rho c_p u)_\mathrm{conv}^* = (\rho c_p u)_\mathrm{conv}
    - L\,\rho_g^\st\,\frac{u_o}{B_o}\,\frac{\pp R_s}{\pp T}.
\end{equation}
Аналогично, часть $Q_\mathrm{deg}$, содержащая $\pp p/\pp t$
и $\pp p/\pp r$, поглощается в модифицированные
$Q_\mathrm{ad}^*$ и $Q_\mathrm{JT}^*$.
Перегруппированное уравнение энергии:
\begin{equation}\label{eq:energy_regrouped}
  (\rho c_p)_\eff^*\,\frac{\pp T}{\pp t}
  + (\rho c_p u)_\mathrm{conv}^*\,\frac{\pp T}{\pp r}
  = \frac{1}{r}\,\frac{\pp}{\pp r}\!\left(
      r\,\lambda_\eff\,\frac{\pp T}{\pp r}
    \right)
  + Q_\mathrm{JT}^* + Q_\mathrm{ad}^*,
\end{equation}
Явно:
\begin{equation}\label{eq:QJT_star}
  Q_\mathrm{JT}^*
  = Q_\mathrm{JT}
    + L\,\rho_g^\st\,\frac{u_o}{B_o}\,
      \frac{\pp R_s}{\pp p}\,\frac{\pp p}{\pp r},
\end{equation}
\begin{equation}\label{eq:Qad_star}
  Q_\mathrm{ad}^*
  = Q_\mathrm{ad}
    + L\,\rho_g^\st\,\phi\,\xi_o\,
      \frac{\pp R_s}{\pp p}\,\frac{\pp p}{\pp t}.
\end{equation}
Таким образом, $Q_\mathrm{JT}^*$ и $Q_\mathrm{ad}^*$ зависят
только от $p$, $T$, $S_g$ и производных давления (не от $\pp T$).


\subsection{Граничные и начальные условия}

\begin{equation}\label{eq:bc_p}
  p(r_w, t) = p_{wf}(t), \qquad
  p(R_\ext, t) = p_0.
\end{equation}
\begin{equation}\label{eq:bc_T}
  T(R_\ext, t) = T_0, \qquad
  \frac{\pp T}{\pp r}\bigg|_{r = r_w} = 0.
\end{equation}
\begin{equation}\label{eq:ic}
  p(r, 0) = p_0, \qquad T(r, 0) = T_0, \qquad S_g(r, 0) = 0.
\end{equation}

При $t = 0$ давление и температура однородны, газ полностью
растворён ($p_0 > p_b(T_0)$). При снижении $p_{wf} < p_b$
вблизи скважины возникает зона разгазирования
$r \in [r_w, r_b(t)]$, где $p < p_b(T)$.


%=============================================================================
\section{Определение интегральных моментов}\label{sec:moments_def}
%=============================================================================

\subsection{Зона разгазирования}

Радиус зоны разгазирования $r_b(t)$ определяется из:
\begin{equation}\label{eq:rb}
  p(r_b(t), t) = p_b(T(r_b(t), t)).
\end{equation}
При $r > r_b(t)$: $p > p_b(T)$, $S_g = 0$, фильтрация однофазная.

\subsection{Моменты давления}

\begin{equation}\label{eq:Mn}
  M_n(t) = \int_{r_w}^{R_\ext} r^{2n+1}\,p(r,t)\,\dd r,
  \qquad n = 0, 1, 2, \ldots
\end{equation}

\paragraph{Нулевой момент ($n = 0$):}
\begin{equation}\label{eq:M0}
  M_0(t) = \int_{r_w}^{R_\ext} r\,p(r,t)\,\dd r.
\end{equation}

\paragraph{Первый момент ($n = 1$):}
\begin{equation}\label{eq:M1}
  M_1(t) = \int_{r_w}^{R_\ext} r^3\,p(r,t)\,\dd r.
\end{equation}

\subsection{Момент насыщенности}

Газонасыщенность имеет компактный носитель $[r_w, r_b(t)]$:
\begin{equation}\label{eq:Sn}
  \mathcal{S}_n(t) = \int_{r_w}^{r_b(t)} r^{2n+1}\,S_g(r,t)\,\dd r.
\end{equation}

$\mathcal{S}_0$ пропорционален объёму свободного газа
в зоне разгазирования.

\subsection{Момент температуры}

Температурное поле определено на всей области $[r_w, R_\ext]$.
Нулевой момент:
\begin{equation}\label{eq:Tn}
  \mathcal{T}_0(t) = \int_{r_w}^{R_\ext} r\,T(r,t)\,\dd r.
\end{equation}

Физический смысл: $2\pi h\,\mathcal{T}_0$ пропорционален
полной тепловой энергии (отклонению от $T_0$) в пласте.


%=============================================================================
\section{Интегральные соотношения}\label{sec:integral_rel}
%=============================================================================

\subsection{Нулевой момент уравнения давления}

Умножим~\eqref{eq:p_eq} на $r$ и проинтегрируем по $r$
от $r_w$ до $R_\ext$. Используя $\mathcal{D}$
из~\eqref{eq:D_operator}:
\begin{multline}\label{eq:moment0}
  \boxed{
    \int_{r_w}^{R_\ext} r\,C_p\,\frac{\pp p}{\pp t}\,\dd r
    + \int_{r_w}^{R_\ext} r\,C_T\,\frac{\pp T}{\pp t}\,\dd r
  } \\
  \boxed{
    = \left[r\,\Lambda\,\frac{\pp p}{\pp r}\right]_{r_w}^{R_\ext}
    + \int_{r_w}^{R_\ext}
      r\,\frac{\lambda_o}{B_o}\!\left(
        \frac{\pp R_s}{\pp p}\,\frac{\pp p}{\pp r}
        + \frac{\pp R_s}{\pp T}\,\frac{\pp T}{\pp r}
      \right)\frac{\pp p}{\pp r}\,\dd r.
  }
\end{multline}

По сравнению с изотермическим случаем появились:
\begin{itemize}[nosep]
  \item член $\int r\,C_T\,(\pp T/\pp t)\,\dd r$ в левой части;
  \item конвективный перенос газа за счёт градиента температуры
        $(\pp R_s/\pp T)\,(\pp T/\pp r)\,(\pp p/\pp r)$ в квадратичном
        интеграле.
\end{itemize}

\paragraph{Вычисление квадратичного интеграла.}
В зоне разгазирования координаты $\eta$ и $\zeta$
связаны линейно: $\zeta = 1 - (L_b/L_0)\,\eta$.
Поэтому произведение $(\pp p/\pp r)\,(\pp T/\pp r)$
при замороженных PVT-свойствах выражается через степенные
функции $\eta$:
\[
  \frac{\pp T}{\pp r}\,\frac{\pp p}{\pp r}
  = \frac{2\sigma\,\Delta T_w\,\Delta p_1\,\gamma}
    {r^2\,L_0\,L_b}\,
    \zeta^{\sigma-1}(1-\zeta^\sigma)\,\eta^{\gamma-1},
\]
где $\zeta(\eta) = 1-(L_b/L_0)\eta$.
При $L_b \ll L_0$ (типично: $L_b/L_0 \sim 0{,}01$--$0{,}1$)
$\zeta \approx 1$ в зоне разгазирования, и
интеграл упрощается.
В однофазной зоне ($\eta$ не определена, но $S_g = 0$)
вклад $\pp R_s/\pp p = 0$ при $p > p_b^*$, и
квадратичный интеграл обнуляется.

\subsection{Первый момент уравнения давления}

Умножим~\eqref{eq:p_eq} на $r^3$ и проинтегрируем:
\begin{multline}\label{eq:moment1}
  \int_{r_w}^{R_\ext} r^3\,C_p\,\frac{\pp p}{\pp t}\,\dd r
  + \int_{r_w}^{R_\ext} r^3\,C_T\,\frac{\pp T}{\pp t}\,\dd r \\
  = \left[r^3\,\Lambda\,\frac{\pp p}{\pp r}\right]_{r_w}^{R_\ext}
    - 2\int_{r_w}^{R_\ext} r^2\,\Lambda\,\frac{\pp p}{\pp r}\,\dd r
  + \int_{r_w}^{R_\ext}
    r^3\,\frac{\lambda_o}{B_o}\!\left(
      \frac{\pp R_s}{\pp p}\,\frac{\pp p}{\pp r}
      + \frac{\pp R_s}{\pp T}\,\frac{\pp T}{\pp r}
    \right)\frac{\pp p}{\pp r}\,\dd r.
\end{multline}

\subsection{Нулевой момент уравнения насыщенности}

Умножим~\eqref{eq:Sg_eq2} на $r$ и проинтегрируем от $r_w$ до $r_b$:
\begin{multline}\label{eq:Sg_moment0}
  \frac{\dd\mathcal{S}_0}{\dd t}
  = \int_{r_w}^{r_b} r\,S_o\,(c_r+c_o^\sat)\,
      \frac{\pp p}{\pp t}\,\dd r
    - \int_{r_w}^{r_b} r\,S_o\,\beta_o^\sat\,
      \frac{\pp T}{\pp t}\,\dd r \\
  - \frac{B_o}{\phi}\!\left[
      \frac{r\,\lambda_o}{B_o}\,\frac{\pp p}{\pp r}
    \right]_{r_w}^{r_b}.
\end{multline}

Член с $\dot{r}_b$ обращается в нуль ($S_g(r_b,t) = 0$,
правило Лейбница). Граничные потоки:
\begin{equation}\label{eq:Sg_flux}
  \left[\frac{r\,\lambda_o}{B_o}\,\frac{\pp p}{\pp r}\right]_{r_w}^{r_b}
  = \frac{r_b\,\lambda_o(r_b)}{B_o(p_b,T_b)}\,p'_b
    - \frac{r_w\,\lambda_o(r_w)}{B_o(p_{wf},T_w)}\,p'_w,
\end{equation}
где $T_b = T(r_b,t)$, $T_w = T(r_w,t)$.

По сравнению с изотермическим случаем появился
интеграл $-\int r\,S_o\,\beta_o^\sat\,(\pp T/\pp t)\,\dd r$.


\subsection{Нулевой момент уравнения энергии}

Умножим~\eqref{eq:energy_regrouped} на $r$ и проинтегрируем
от $r_w$ до $R_\ext$:
\begin{multline}\label{eq:T_moment0}
  \int_{r_w}^{R_\ext} r\,(\rho c_p)_\eff^*\,
    \frac{\pp T}{\pp t}\,\dd r
  + \int_{r_w}^{R_\ext} r\,(\rho c_p u)_\mathrm{conv}^*\,
    \frac{\pp T}{\pp r}\,\dd r \\
  = \left[r\,\lambda_\eff\,\frac{\pp T}{\pp r}
    \right]_{r_w}^{R_\ext}
  + \int_{r_w}^{R_\ext} r\,(Q_\mathrm{JT}^* + Q_\mathrm{ad}^*)\,\dd r.
\end{multline}

Дивергентная часть вычисляется точно:
\begin{equation}\label{eq:T_flux}
  \left[r\,\lambda_\eff\,\frac{\pp T}{\pp r}
    \right]_{r_w}^{R_\ext}
  = R_\ext\,\lambda_\eff\,T'_\ext
    - \underbrace{r_w\,\lambda_\eff\,T'_w}_{= 0},
\end{equation}
поскольку граничное условие~\eqref{eq:bc_T} даёт $T'_w = 0$.


%=============================================================================
\section{Параметризация профилей}\label{sec:profiles}
%=============================================================================

\subsection{Структура области}

Область $[r_w, R_\ext]$ делится на две зоны:
\begin{itemize}[nosep]
  \item \textbf{Зона разгазирования} $r \in [r_w, r_b(t)]$:
        $p < p_b(T)$, $S_g > 0$, двухфазная фильтрация.
  \item \textbf{Однофазная зона} $r \in [r_b(t), R_\ext]$:
        $p \geqslant p_b(T)$, $S_g = 0$, однофазная фильтрация нефти.
\end{itemize}

\subsection{Профиль давления}

\paragraph{Однофазная зона ($r_b \leqslant r \leqslant R_\ext$).}
Квазистационарный логарифмический профиль:
\begin{equation}\label{eq:p_outer}
  p(r,t) = p_b^* + (p_0 - p_b^*)\,
    \frac{\ln(r/r_b)}{\ln(R_\ext/r_b)},
  \qquad r_b \leqslant r \leqslant R_\ext.
\end{equation}

\paragraph{Зона разгазирования ($r_w \leqslant r \leqslant r_b$).}
Степенной профиль:
\begin{equation}\label{eq:p_inner}
  p(r,t) = p_{wf} + (p_b^* - p_{wf})\,\Phi\!\left(
    \frac{\ln(r/r_w)}{\ln(r_b/r_w)}
  \right),
\end{equation}
\begin{equation}\label{eq:Phi}
  \Phi(\eta) = \eta^\gamma, \qquad
  \eta = \frac{\ln(r/r_w)}{\ln(r_b/r_w)},
  \qquad \gamma > 0.
\end{equation}

\paragraph{Непрерывность потока на $r = r_b$.}
При $S_g(r_b) = 0$ подвижности совпадают.
Обозначим $p_b^* = p_b(T(r_b))$ --- давление насыщения
при температуре на фронте разгазирования.
Тогда $\Delta p_1 = p_b^* - p_{wf}$,
$\Delta p_2 = p_0 - p_b^*$, и
\begin{equation}\label{eq:gamma_constraint}
  \gamma = \frac{\Delta p_2\,L_b}{\Delta p_1\,L_e}.
\end{equation}
Величина $p_b^* = p_b(T(r_b))$ зависит от $\Delta T_w$
через профиль~\eqref{eq:T_profile}:
\begin{equation}\label{eq:pb_star}
  p_b^* = p_b\!\left(T_0 - \Delta T_w\,\Psi(\zeta_b)\right),
  \qquad
  \zeta_b = \frac{\ln(R_\ext/r_b)}{L_0} = \frac{L_e}{L_0}.
\end{equation}
При $\Delta T_w = 0$ (начальная стадия):
$p_b^* = p_b(T_0)$ --- фиксированное значение.

Производные давления:
\begin{equation}\label{eq:dp_inner}
  \frac{\pp p}{\pp r}\bigg|_\text{inner}
  = \frac{\Delta p_1\,\gamma}{r\,L_b}\,\eta^{\gamma-1},
\end{equation}
\begin{equation}\label{eq:dp_outer}
  \frac{\pp p}{\pp r}\bigg|_\text{outer}
  = \frac{\Delta p_2}{r\,L_e}.
\end{equation}

\paragraph{Обозначения.}
$L_b = \ln(r_b/r_w)$, $L_e = \ln(R_\ext/r_b)$,
$\Delta p_1 = p_b^* - p_{wf}$, $\Delta p_2 = p_0 - p_b^*$
(здесь и далее $p_b^* = p_b(T(r_b))$, см.~\eqref{eq:pb_star}).


\subsection{Профиль насыщенности}

Газонасыщенность имеет компактный носитель $[r_w, r_b(t)]$:
\begin{equation}\label{eq:Sg_profile}
  S_g(r,t) = S_{g,w}(t)\!\left[
    1 - \left(\frac{\ln(r/r_w)}{\ln(r_b/r_w)}\right)^{\!\beta}
  \right],
  \qquad r_w \leqslant r \leqslant r_b,
\end{equation}
$S_g(r_w) = S_{g,w}$, $S_g(r_b) = 0$, $\beta > 0$.


\subsection{Профиль температуры}\label{ssec:T_profile}

Температурное поле определено на $[r_w, R_\ext]$ и удовлетворяет:
\begin{itemize}[nosep]
  \item $T(R_\ext, t) = T_0$ (контур питания);
  \item $T'(r_w, t) = 0$ (адиабатическая стенка).
\end{itemize}

Профиль параметризуем функцией формы
с амплитудным параметром $\Delta T_w(t) = T_0 - T_w(t)$,
где $T_w(t) = T(r_w, t)$ --- температура на забое:
\begin{equation}\label{eq:T_profile}
  T(r,t) = T_0 - \Delta T_w(t)\,\Psi(\zeta),
  \qquad
  \zeta = \frac{\ln(R_\ext/r)}{\ln(R_\ext/r_w)} = \frac{\ln(R_\ext/r)}{L_0}.
\end{equation}

Функция формы $\Psi(\zeta)$ должна удовлетворять:
$\Psi(0) = 0$ (контур питания),
$\Psi(1) = 1$ (забой),
$\Psi'(1) = 0$ (адиабатическая стенка: $\pp T/\pp r|_{r_w} = 0$).

Простейший профиль с этими свойствами:
\begin{equation}\label{eq:Psi}
  \boxed{
    \Psi(\zeta) = \zeta^\sigma(2 - \zeta^\sigma),
    \qquad \sigma > 0.
  }
\end{equation}

Проверка:
\begin{itemize}[nosep]
  \item $\Psi(0) = 0$. \checkmark
  \item $\Psi(1) = 1\cdot(2-1) = 1$. \checkmark
  \item $\Psi'(\zeta) = 2\sigma\zeta^{\sigma-1}(1 - \zeta^\sigma)$,
        поэтому $\Psi'(1) = 2\sigma\cdot 1\cdot 0 = 0$. \checkmark
\end{itemize}

Производная температуры:
\begin{equation}\label{eq:dT_dr}
  \frac{\pp T}{\pp r}
  = \frac{\Delta T_w}{r\,L_0}\,\Psi'(\zeta)
  = \frac{2\sigma\,\Delta T_w}{r\,L_0}\,
    \zeta^{\sigma-1}(1-\zeta^\sigma).
\end{equation}
При $r = r_w$ ($\zeta = 1$): $\pp T/\pp r = 0$. \checkmark
Максимум $|\pp T/\pp r|$ достигается внутри области.

\paragraph{Интегралы.}
Моменты профиля выражаются через:
\begin{equation}\label{eq:IPsi}
  J^{(m)} = \int_0^1 \zeta^m\,(2 - \zeta^\sigma)\,\dd\zeta
  = \frac{2}{m+1} - \frac{1}{m+\sigma+1}.
\end{equation}

\paragraph{Выбор $\sigma$.}
При преобладании конвекции ($\mathrm{Pe} \gg 1$)
нефтяной поток подводит невозмущённую
температуру $T_0$ от контура, а охлаждение
генерируется локально (JT-эффект, дегазация).
Возмущение $\Delta T$ концентрируется
вблизи $r_w$. При $\sigma = 1$: $\Psi = \zeta(2-\zeta) = 1-(1-\zeta)^2$ ---
параболический профиль, основное изменение
вблизи забоя. При $\sigma > 1$ профиль
ещё сильнее локализован у $r_w$.

$\sigma$ можно зафиксировать ($\sigma = 1$) или определять
из первого момента уравнения энергии (добавив четвёртое ODE).
Начинаем с фиксированного $\sigma = 1$.


%=============================================================================
\section{Система ODE}\label{sec:ode}
%=============================================================================

\subsection{Свободные параметры}

Профили содержат зависящие от времени параметры:
\begin{center}
\begin{tabular}{cl}
  $r_b(t)$       & радиус зоны разгазирования, \\
  $\gamma(t)$    & показатель формы давления
                    ($\gamma = \gamma(r_b)$
                    из~\eqref{eq:gamma_constraint}), \\
  $S_{g,w}(t)$   & газонасыщенность на забое, \\
  $\Delta T_w(t)$ & отклонение температуры на забое от $T_0$.
\end{tabular}
\end{center}

Показатели формы $\beta$ и $\sigma$ фиксированы.
$\gamma$ --- не независимый: $\gamma = \gamma(r_b)$.
Остаются \textbf{три свободных параметра}:
$r_b(t)$, $S_{g,w}(t)$, $\Delta T_w(t)$.

\subsection{Уравнения для параметров}

Подставляя профили~\eqref{eq:p_outer}--\eqref{eq:T_profile}
в три интегральных соотношения~\eqref{eq:moment0},
\eqref{eq:Sg_moment0}, \eqref{eq:T_moment0},
получаем систему трёх ODE:

\paragraph{1. Из нулевого момента давления~\eqref{eq:moment0}:}
\begin{equation}\label{eq:ode1}
  \mathcal{A}_{11}\,\dot{r}_b
  + \mathcal{A}_{12}\,\dot{S}_{g,w}
  + \mathcal{A}_{13}\,\dot{\Delta T}_w
  = \mathcal{F}_1.
\end{equation}

\paragraph{2. Из нулевого момента насыщенности~\eqref{eq:Sg_moment0}:}
\begin{equation}\label{eq:ode2}
  \mathcal{A}_{21}\,\dot{r}_b
  + \mathcal{A}_{22}\,\dot{S}_{g,w}
  + \mathcal{A}_{23}\,\dot{\Delta T}_w
  = \mathcal{F}_2.
\end{equation}

\paragraph{3. Из нулевого момента энергии~\eqref{eq:T_moment0}:}
\begin{equation}\label{eq:ode3}
  \mathcal{A}_{31}\,\dot{r}_b
  + \mathcal{A}_{32}\,\dot{S}_{g,w}
  + \mathcal{A}_{33}\,\dot{\Delta T}_w
  = \mathcal{F}_3.
\end{equation}

Матричная форма:
\begin{equation}\label{eq:ode_system}
  \boxed{
    \begin{pmatrix}
      \mathcal{A}_{11} & \mathcal{A}_{12} & \mathcal{A}_{13} \\
      \mathcal{A}_{21} & \mathcal{A}_{22} & \mathcal{A}_{23} \\
      \mathcal{A}_{31} & \mathcal{A}_{32} & \mathcal{A}_{33}
    \end{pmatrix}
    \begin{pmatrix}
      \dot{r}_b \\ \dot{S}_{g,w} \\ \dot{\Delta T}_w
    \end{pmatrix}
    =
    \begin{pmatrix}
      \mathcal{F}_1 \\ \mathcal{F}_2 \\ \mathcal{F}_3
    \end{pmatrix}.
  }
\end{equation}

\subsection{Структура коэффициентов}

Коэффициенты $\mathcal{A}_{ij}$ и $\mathcal{F}_i$ вычисляются
подстановкой профилей~\eqref{eq:p_outer}--\eqref{eq:T_profile}
в интегралы и дифференцированием по параметрам.

\paragraph{Связь параметров с производными профилей.}
Профили $p(r,t)$ и $T(r,t)$ зависят от $t$ через
$r_b(t)$, $\Delta T_w(t)$ (и $p_{wf}(t)$, задаваемое).
Поэтому $\pp p/\pp t$ и $\pp T/\pp t$ выражаются
через $\dot{r}_b$ и $\dot{\Delta T}_w$:
\begin{equation}\label{eq:dp_dt}
  \frac{\pp p}{\pp t}
  = \frac{\pp p}{\pp r_b}\,\dot{r}_b
    + \frac{\pp p}{\pp p_{wf}}\,\dot{p}_{wf},
\end{equation}
\begin{equation}\label{eq:dT_dt}
  \frac{\pp T}{\pp t}
  = \frac{\pp T}{\pp r_b}\,\dot{r}_b
    + \frac{\pp T}{\pp \Delta T_w}\,\dot{\Delta T}_w.
\end{equation}
Профиль $S_g$ зависит от $r_b$ и $S_{g,w}$, поэтому
$\pp S_g/\pp t$ содержит $\dot{r}_b$ и $\dot{S}_{g,w}$.

Подставляя~\eqref{eq:dp_dt}--\eqref{eq:dT_dt}
в~\eqref{eq:moment0}--\eqref{eq:T_moment0}, собираем
коэффициенты при $\dot{r}_b$, $\dot{S}_{g,w}$,
$\dot{\Delta T}_w$ --- это элементы матрицы
$\mathcal{A}_{ij}$. Члены, не содержащие производных
по времени, образуют правую часть $\mathcal{F}_i$.

\paragraph{Температурный вклад в уравнение давления.}
Элемент $\mathcal{A}_{13}$ содержит:
\[
  \mathcal{A}_{13} \sim \int_{r_w}^{R_\ext} r\,C_T\,
    \frac{\pp T}{\pp \Delta T_w}\,\dd r.
\]
Из~\eqref{eq:T_profile}: $\pp T/\pp \Delta T_w = -\Psi(\zeta)$.
При $C_T < 0$ (типичный случай: $\pp R_s/\pp T < 0$)
этот вклад положителен: охлаждение увеличивает
давление насыщения и <<тормозит>> расширение зоны разгазирования.

\paragraph{Температурный вклад в уравнение насыщенности.}
Элемент $\mathcal{A}_{23}$ содержит:
\[
  \mathcal{A}_{23} \sim -\int_{r_w}^{r_b} r\,S_o\,\beta_o^\sat\,
    \frac{\pp T}{\pp \Delta T_w}\,\dd r.
\]


\subsection{Явный вид коэффициентов}

\paragraph{Двухзонное замораживание PVT-свойств.}
Для вычисления интегралов PVT-свойства замораживаются
раздельно в каждой зоне:
\begin{itemize}[nosep]
  \item \textbf{Зона разгазирования} ($r_w \leqslant r \leqslant r_b$):
        \[
          \bar{p}_1 = \frac{p_{wf} + p_b^*}{2}, \qquad
          \bar{T}_1 = T_0 - \Delta T_w\,\bar{\Psi}_1,
          \qquad
          \bar{S}_{g,1} = \frac{S_{g,w}}{2},
        \]
        где
        \[
          \bar{\Psi}_1
          = \frac{1}{L_b}\int_{\zeta_b}^{1}\Psi(\zeta)\,\dd\zeta.
        \]
        При $L_b \ll L_0$ ($\zeta_b \to 1$):
        $\bar{\Psi}_1 \to 1$, т.\,е.\ $\bar{T}_1 \approx T_0 - \Delta T_w$
        (вся зона разгазирования лежит вблизи забоя).
  \item \textbf{Однофазная зона} ($r_b \leqslant r \leqslant R_\ext$):
        \[
          \bar{p}_2 = \frac{p_b^* + p_0}{2}, \qquad
          \bar{T}_2 \approx T_0, \qquad
          S_g = 0.
        \]
\end{itemize}

Для компактности введём обозначения:
\begin{align}
  I_\Phi^{(n)} &= \int_0^1 \eta^{n\gamma}\,\dd\eta
    = \frac{1}{n\gamma + 1}, \label{eq:IPhi} \\
  I_S^{(n)} &= \int_0^1 (1-\eta^\beta)\,\eta^{n\gamma}\,\dd\eta
    = I_\Phi^{(n)} - \frac{1}{n\gamma + \beta + 1}.
    \label{eq:IS}
\end{align}
Интегралы профиля температуры $J^{(m)}$
определены в~\eqref{eq:IPsi}.

\paragraph{Расщепление матрицы.}
Профили $p(r)$ и $T(r)$ не зависят от $S_{g,w}$
(давление определяется через $r_b$, $p_{wf}$;
температура --- через $r_b$, $\Delta T_w$).
Поэтому $\pp p/\pp t$ и $\pp T/\pp t$ не содержат
$\dot{S}_{g,w}$, и
\begin{equation}\label{eq:matrix_structure}
  \mathcal{A}_{12} = 0, \qquad \mathcal{A}_{32} = 0.
\end{equation}
Матрица системы~\eqref{eq:ode_system} имеет вид:
\[
  \begin{pmatrix}
    \mathcal{A}_{11} & 0 & \mathcal{A}_{13} \\
    \mathcal{A}_{21} & \mathcal{A}_{22} & \mathcal{A}_{23} \\
    \mathcal{A}_{31} & 0 & \mathcal{A}_{33}
  \end{pmatrix}.
\]
Из первого и третьего уравнений определяются $\dot{r}_b$
и $\dot{\Delta T}_w$ (подсистема $2\times 2$), после чего
из второго уравнения находится $\dot{S}_{g,w}$.

\paragraph{Замечание.}
Расщепление касается только производных $\dot{S}_{g,w}$.
Величина $S_{g,w}$ входит в коэффициенты $\mathcal{A}_{ij}$
и $\mathcal{F}_i$ через PVT-свойства (оцениваемые
при замороженных значениях $\bar{S}_g$), поэтому
система остаётся связанной.

Явные формулы для всех $\mathcal{A}_{ij}$ и $\mathcal{F}_i$
через эти интегралы приводятся при конкретной реализации
(вычислительная записка).


%=============================================================================
\section{Начальная стадия: автомодельная асимптотика при $t \to 0$}
\label{sec:initial_stage}
%=============================================================================

При $t = 0$ давление скачком снижается до $p_{wf} < p_b$.
Формально $r_b(0) = r_w$, $S_{g,w}(0) = 0$, $\Delta T_w(0) = 0$.
При $t = 0^+$ $S_{g,w}$ скачком принимает конечное
значение, но $L_b = 0$, $\gamma = 0$, и матрица
системы~\eqref{eq:ode_system} вырождается.

\subsection{Физическая картина}

При $t < 0$ давление $p = p_0 > p_b$, температура $T = T_0$:
газ растворён, фильтрация однофазная.
При $t = 0$ мгновенно $p(r_w) = p_{wf} < p_b$.

На малых временах $t \ll r_w^2/\chi_0$ задача описывается
однофазной диффузией давления. Температура практически
не меняется: характерное время теплового отклика
$\tau_T \sim r_w^2/a$ ($a = \lambda_\eff/(\rho c_p)_\eff$
--- температуропроводность), и при типичных параметрах
$a \sim 10^{-7}$~м$^2$/с, $r_w = 0{,}1$~м:
$\tau_T \sim 10^5$~с, т.\,е.\
$\tau_T/\tau_\mathrm{rel} \sim 10^7$.

Таким образом, на стадии зарождения зоны разгазирования
температурное поле <<заморожено>>: $T \approx T_0$,
$\Delta T_w \approx 0$. Начальные условия для $r_b$
и $S_{g,w}$ определяются из изотермической задачи.

\subsection{Локальная задача диффузии}

При $\delta = r_b - r_w \ll r_w$:
\begin{equation}\label{eq:diffusion}
  \frac{\pp p}{\pp t} = \chi_0\,\frac{\pp^2 p}{\pp \xi^2},
  \qquad
  \chi_0 = \frac{\Lambda_0}{C_p^{(0)}},
  \qquad
  \xi = r - r_w,
\end{equation}
где $\Lambda_0 = k/(B_g\,\mu_o)$ при $S_g = 0$,
$C_p^{(0)}$ --- коэффициент~\eqref{eq:Cp}
при $S_g = 0$, $T = T_0$.
Член $C_T\,\dot{T}$ в~\eqref{eq:p_eq} обнуляется не потому,
что $C_T = 0$, а потому, что на этой стадии $\dot{T} \approx 0$
($\tau_T \gg \tau_\mathrm{rel}$).

Решение:
\begin{equation}\label{eq:p_erfc}
  p(\xi, t) = p_{wf} + (p_0 - p_{wf})\,
    \operatorname{erf}\!\left(\frac{\xi}{2\sqrt{\chi_0\,t}}\right).
\end{equation}

\subsection{Скорость роста границы разгазирования}

Условие $p(r_b) = p_b$ даёт:
\begin{equation}\label{eq:delta_sqrt}
  \boxed{
    \delta(t) = r_b(t) - r_w = 2\,\alpha_*\,\sqrt{\chi_0\,t},
    \qquad
    \alpha_* = \operatorname{erf}^{-1}\!\left(
      \frac{p_b - p_{wf}}{p_0 - p_{wf}}
    \right).
  }
\end{equation}

В отличие от задачи Стефана, $r_b$ --- кинематически
определённая изобара на гладком erfc-профиле.

\subsection{Асимптотика параметров профилей}

При $\delta \ll r_w$:
\begin{equation}\label{eq:Lb_asymp}
  L_b \approx \frac{\delta}{r_w},
  \qquad
  L_e \to L_0 = \ln\!\frac{R_\ext}{r_w}.
\end{equation}
\begin{equation}\label{eq:gamma_asymp}
  \gamma(t)
  \approx \frac{\Delta p_2}{\Delta p_1\,L_0}\,
    \frac{2\,\alpha_*\sqrt{\chi_0\,t}}{r_w} \to 0.
\end{equation}

\subsection{Начальная газонасыщенность}

Квазиравновесие (flash):
\begin{equation}\label{eq:Sgw_equil}
  \boxed{
    S_{g,w}^{\mathrm{eq}}
    = \frac{[R_s(p_b, T_0) - R_s(p_{wf}, T_0)]\,
      B_g(p_{wf}, T_0)}{B_o(p_{wf}, T_0)}.
  }
\end{equation}

$S_{g,w}^{\mathrm{eq}}$ устанавливается при $t = 0^+$
(локальное равновесие Black~Oil, $T = T_0$).
Формула~\eqref{eq:Sgw_equil} --- линеаризация при $S_g \ll 1$;
точное выражение содержит $B_o + \Delta R_s\,B_g$
в знаменателе вместо~$B_o$.
Конвективный перенос газа потоком нефти
при $S_g < S_{g,\mathrm{cr}}$ снижает $S_{g,w}$,
делая~\eqref{eq:Sgw_equil} верхней оценкой.

\subsection{Начальная температура}

При $t = t_\varepsilon$: $\Delta T_w(t_\varepsilon) \approx 0$.
Температура ещё не успела измениться.

Первое ненулевое приближение: охлаждение определяется
балансом эффекта Джоуля--Томсона и теплоты дегазации.
Интегрируя~\eqref{eq:energy} по малому объёму
$r_w < r < r_w + \delta$ при $T'_w = 0$:
\begin{equation}\label{eq:dTw_init}
  (\rho c_p)_\eff\,\frac{\dd T_w}{\dd t}\bigg|_{t=0^+}
  \approx Q_\mathrm{JT}(r_w) + Q_\mathrm{ad}(r_w) + Q_\mathrm{deg}(r_w).
\end{equation}
Оба источника конечны при $t = 0^+$ (давление и его
градиент определены из erfc-решения), поэтому
$\Delta T_w(t_\varepsilon) = O(t_\varepsilon) \approx 0$
при $t_\varepsilon \sim 10^{-6}\text{--}10^{-4}$~с.

\subsection{Сводка: начальные данные для ODE-системы}

Для запуска интегрирования системы~\eqref{eq:ode_system}
при $t = t_\varepsilon > 0$:
\begin{equation}\label{eq:ic_regularized}
  \boxed{
    r_b(t_\varepsilon) = r_w + 2\,\alpha_*\sqrt{\chi_0\,t_\varepsilon},
    \quad
    S_{g,w}(t_\varepsilon) = S_{g,w}^{\mathrm{eq}},
    \quad
    \Delta T_w(t_\varepsilon) = 0.
  }
\end{equation}

$t_\varepsilon$ выбирается из условия $\delta/r_w \sim 10^{-2}\text{--}10^{-1}$,
что при $\chi_0 \sim 0{,}5$~м$^2$/с, $\alpha_* \approx 0{,}5$,
$r_w = 0{,}1$~м даёт $t_\varepsilon \sim 10^{-6}\text{--}10^{-4}$~с.
На этих временах $\Delta T_w \approx 0$: тепловой отклик пренебрежимо мал.


%=============================================================================
\section{Обсуждение}\label{sec:discussion}
%=============================================================================

\subsection{Преимущества подхода}

\begin{enumerate}[nosep]
  \item Три PDE редуцированы к системе трёх ODE ---
        вычислительная стоимость на порядки ниже.
  \item Подвижная граница $r_b(t)$ --- естественный параметр.
  \item Температурное поле включено в систему: обратная связь
        $T \to R_s(p,T)$ учитывается автоматически на каждом
        временном шаге.
  \item Интегральные характеристики (дебит, объём газа,
        средняя температура) воспроизводятся по построению.
\end{enumerate}

\subsection{Сложность решения ODE-системы}

Формально система~\eqref{eq:ode_system} состоит из трёх ODE,
что на порядки дешевле исходных PDE. Однако реальная сложность
определяется не размерностью, а структурой правых частей:

\begin{enumerate}[nosep]
  \item \textbf{Вычисление коэффициентов.}
        На каждом временном шаге необходимо вычислить 9 элементов
        матрицы $\mathcal{A}_{ij}$ (из которых 2 нулевых)
        и 3~компоненты $\mathcal{F}_i$. Каждый элемент ---
        линейная комбинация интегралов $I_\Phi^{(n)}$,
        $I_\Psi^{(n)}$, $I_S^{(n)}$, которые при замораживании
        PVT-свойств вычисляются аналитически.
        Стоимость одного шага --- $O(1)$.

  \item \textbf{Конвективный интеграл в уравнении энергии.}
        Подставляя $u_\alpha = -\lambda_\alpha\,\pp p/\pp r$
        и~\eqref{eq:dT_dr}:
        \[
          \int_{r_w}^{R_\ext} r\,(\rho c_p u)_\mathrm{conv}^*\,
            \frac{\pp T}{\pp r}\,\dd r
          = -\int_{r_w}^{R_\ext}
            \frac{(\rho c_p \lambda)_\mathrm{conv}^*}{r}\,
            \frac{\pp p}{\pp r}\,
            \Psi'(\zeta)\,\frac{\Delta T_w}{L_0}\,\dd r,
        \]
        где $(\rho c_p \lambda)_\mathrm{conv}^*
        = \sum_\alpha \rho_\alpha c_{p,\alpha}\,\lambda_\alpha
          - L\,\rho_g^\st\,(\lambda_o/B_o)\,\pp R_s/\pp T$.
        При замораживании коэффициентов и подстановке
        $\pp p/\pp r$ из~\eqref{eq:dp_inner}--\eqref{eq:dp_outer}
        интеграл сводится к комбинации $I_\Phi^{(n)}$
        и $J^{(m)}$ через связь $\zeta(\eta)$.

  \item \textbf{Квадратичный член в уравнении давления.}
        Интеграл $\int r\,(\lambda_o/B_o)\,\nabla R_s\cdot\nabla p\,
        |\nabla p|\,\dd r$ содержит $(\pp p/\pp r)^2$ и
        $(\pp T/\pp r)\,(\pp p/\pp r)$. Второе слагаемое ---
        произведение производных \emph{разных} профилей,
        определённых в \emph{разных} координатах
        ($\eta$ для $p$, $\zeta$ для $T$), что не сводится
        к стандартным $I$-интегралам.

  \item \textbf{Жёсткость.}
        Система потенциально жёсткая: масштаб $\dot{r}_b$
        определяется гидродинамикой ($\chi_0 \sim 0{,}5$~м$^2$/с),
        а $\dot{\Delta T}_w$ --- температуропроводностью
        ($a \sim 10^{-7}$~м$^2$/с). Отношение масштабов
        $\chi_0/a \sim 10^6$--$10^7$, что типично
        для жёстких систем.
        Конвективный перенос ($\mathrm{Pe} \gg 1$) смягчает
        жёсткость: эффективная температуропроводность
        $a_\eff \sim uL \sim 10^{-4}$~м$^2$/с
        ($u \sim 10^{-3}$~м/с, $L \sim r_w$), что
        снижает отношение до $\chi_0/a_\eff \sim 10^3$--$10^4$.
        Система остаётся жёсткой; необходим неявный
        интегратор (BDF порядка 2--5).

  \item \textbf{Расщепление.}
        Структура~\eqref{eq:matrix_structure} позволяет
        решать подсистему $2\times 2$ на $(r_b, \Delta T_w)$
        и скалярное уравнение на $S_{g,w}$ последовательно,
        что упрощает обращение матрицы.
\end{enumerate}

\paragraph{Обезразмеривание.}
Элементы $\mathcal{A}_{11}$ масштабируются как $C_p$~[Па$^{-1}$],
а $\mathcal{A}_{33}$ --- как $(\rho c_p)_\eff^*$~[Дж/(м$^3$\,К)].
Без обезразмеривания обусловленность матрицы $\sim 10^{15}$.
Рекомендуется ввести характерные масштабы
$[p] = p_0 - p_{wf}$, $[T] = 10$~К, $[r] = r_w$, $[t] = r_w^2/\chi_0$
и перейти к безразмерным переменным $\tilde{r}_b$,
$\tilde{\Delta T}_w$, $\tilde{S}_{g,w}$.

\paragraph{Итого.}
При замораживании PVT-свойств система~\eqref{eq:ode_system}
допускает решение BDF-методом.
Основная вычислительная стоимость --- не в решении ODE,
а в вычислении PVT-свойств при обновлении средних $\bar{p}$,
$\bar{T}$, $\bar{S}_g$. Если PVT заданы корреляциями (Стендинг
и т.\,п.), стоимость одного шага $\sim 10^2$ операций.


\subsection{Ограничения и открытые вопросы}

\paragraph{Решённые вопросы.}
\begin{itemize}[nosep]
  \item Профиль температуры~\eqref{eq:Psi} удовлетворяет
        $\Psi'(1) = 0$, что обеспечивает $T'(r_w) = 0$
        (адиабатическое ГУ на забое).
  \item Зависимость $p_b(T)$ учтена через
        $p_b^* = p_b\bigl(T(r_b)\bigr)$
        в профиле давления и ограничении~\eqref{eq:gamma_constraint}.
  \item Двухзонное замораживание PVT-свойств
        (раздельные средние для зон разгазирования и однофазной)
        снижает ошибку по сравнению с единым средним.
  \item Смешанный интеграл $\int r\,(\pp T/\pp r)\,(\pp p/\pp r)\,\dd r$
        вычисляется аналитически благодаря линейной связи
        $\zeta = 1 - (L_b/L_0)\eta$.
  \item Модифицированные источники $Q_\mathrm{JT}^*$,
        $Q_\mathrm{ad}^*$ выписаны
        явно~\eqref{eq:QJT_star}--\eqref{eq:Qad_star}.
  \item Конвективный интеграл в уравнении энергии
        выражен через замороженные коэффициенты
        и стандартные $I$-интегралы.
\end{itemize}

\paragraph{Открытые вопросы.}
\begin{enumerate}[nosep]
  \item \textbf{Квадратичный член $(\pp p/\pp r)^2$.}
        Интеграл $\int r\,(\lambda_o/B_o)\,(\pp R_s/\pp p)\,
        (\pp p/\pp r)^2\,\dd r$ не сводится к граничным потокам
        и требует замораживания $\lambda_o/B_o$ при среднем давлении.
        Его вклад может быть оценён a~posteriori.

  \item \textbf{Перегруппировка $Q_\mathrm{deg}$.}
        Перегруппированные коэффициенты~\eqref{eq:rhocp_eff_mod},
        \eqref{eq:rhocp_conv_mod} зависят от $p$, $T$, $S_g$,
        что усложняет замораживание. При двухзонной оценке
        эффект смягчается, но строгая оценка погрешности открыта.

  \item \textbf{Выбор $\beta$, $\sigma$.}
        Показатели формы профилей можно зафиксировать или определить
        из старших моментов. Чувствительность подлежит исследованию.

  \item \textbf{Верификация.}
        Сопоставление с численным решением полной PDE-системы
        (метод конечных разностей или OPM Flow)
        необходимо для оценки реальной точности метода моментов.
\end{enumerate}


%=============================================================================
\section*{Литература}
%=============================================================================

\begin{thebibliography}{9}

\bibitem{draft4}
  Двухфазная фильтрация нефти с газом в радиальной геометрии.
  Полная система уравнений с балансом энергии.
  --- Draft4/two\_phase\_energy.tex (рабочий документ).

\bibitem{barenblatt1954}
  Баренблатт~Г.\,И. О некоторых приближённых методах в теории
  одномерной неустановившейся фильтрации жидкости при упругом
  режиме~// Изв. АН СССР, ОТН. --- 1954. --- №\,9. --- С.\,35--49.

\bibitem{BER1984}
  Баренблатт~Г.\,И., Ентов~В.\,М., Рыжик~В.\,М.
  Движение жидкостей и газов в природных пластах. ---
  М.: Недра, 1984. --- 211~с.

\bibitem{mirz2004}
  Мирзаджанзаде~А.\,Х., Хасанов~М.\,М., Бахтизин~Р.\,Н.
  Моделирование процессов нефтегазодобычи. ---
  М.--Ижевск: Институт компьютерных исследований, 2004. ---
  С.\,212--218 (§\,3.7).

\bibitem{pirverdyan1961}
  Пирвердян~А.\,М. Об одном способе оценок приближённых решений
  уравнений нестационарной фильтрации нефти и газа~//
  ПММ. --- 1961. --- Т.\,25, вып.\,4.

\bibitem{entov1965}
  Ентов~В.\,М. Теоремы сравнения для уравнений нестационарной
  фильтрации~// ПММ. --- 1965. --- Т.\,29, вып.\,11.

\end{thebibliography}

\end{document}
